\chapter{Implementação}

\section{Hardware Utilizado}

O hardware da plataforma é composto por duas placas ESP32, um conversor digital-analógico MCP4728, conexões USB para comunicação serial e canais analógicos destinados à reprodução e recepção dos sinais de tensão e corrente.

\begin{table}[H]
  \centering
  \caption{Principais componentes da plataforma}
  \label{tab:componentes}
  \begin{tabular}{ll}
    \toprule
    Componente & Função \\
    \midrule
    ESP32 geradora & Receber amostras e controlar a reprodução dos sinais \\
    MCP4728 & Converter amostras digitais em sinais analógicos \\
    ESP32 receptora & Adquirir sinais analógicos e executar algoritmos de proteção \\
    Computador & Configuração, envio de arquivos e supervisão dos testes \\
    \bottomrule
  \end{tabular}
\end{table}

\section{Software Desenvolvido}

O desenvolvimento do projeto é organizado em scripts e firmwares específicos:

\begin{itemize}
  \item \texttt{comtrade\_generator.m}: gera os arquivos COMTRADE a partir das amostras do Simulink.
  \item \path{run_embedded_generator.py}: valida, interpreta e envia as amostras COMTRADE à ESP32 geradora.
  \item \path{embedded_generator.ino}: firmware responsável pela recepção, bufferização e reprodução dos sinais.
  \item \texttt{run\_embedded\_receiver.py}: executa a interface de controle e supervisão da ESP32 receptora.
  \item \path{embedded_receiver.ino}: firmware responsável pela aquisição das amostras e execução dos algoritmos de proteção.
\end{itemize}

\section{Protocolo de Comunicação}

A comunicação entre o computador e as ESP32 ocorre por interface USB/serial. As amostras são empacotadas em mensagens contendo informações de amplitude, canal e controle de execução. O protocolo deve permitir identificar o início e o fim da transmissão, a separação entre canais de tensão e corrente e a confirmação de recebimento dos blocos de dados.

\section{Bufferização e Reprodução}

A bufferização das amostras é necessária para reduzir falhas de temporização durante a reprodução. O firmware da ESP32 geradora organiza as amostras recebidas em memória e atualiza os canais do DAC conforme a taxa de amostragem definida pelo arquivo COMTRADE. Dessa forma, busca-se preservar a forma de onda original dentro das limitações de resolução, velocidade e faixa de tensão do hardware utilizado.

\section{Aquisição e Processamento}

Na unidade receptora, os sinais analógicos são convertidos em amostras digitais. Em seguida, janelas de amostras são processadas para estimativa das grandezas utilizadas nos algoritmos de proteção. A Transformada de Fourier é aplicada para extração da componente fundamental, permitindo obter magnitude e fase dos sinais de tensão e corrente.

\section{Algoritmos de Proteção}

Os algoritmos de proteção foram estruturados para permitir operação individual ou combinada das funções implementadas. Cada função possui parâmetros próprios de ajuste, configurados de forma semelhante ao processo adotado em relés reais. Essa característica possibilita estudar a influência dos ajustes sobre a atuação do sistema.

Além disso, a arquitetura permite a interação entre funções de proteção. Por exemplo, uma lógica direcional pode ser utilizada em conjunto com uma função de distância, permitindo avaliar estratégias de proteção mais próximas das empregadas em sistemas elétricos reais.
